\chapter{Fit}



\begin{exercise}

Si dimostri che se fittiamo una serie di punti con una constante, la somma dei
residui rispetto al modello di \foreign{best-fit} è nulla.

\end{exercise}



\begin{exercise}

Si dimostri che se fittiamo una serie di punti con una retta, si ha
\begin{align*}
  \sum_i r_i = 0 \quad \text{e} \quad \sum_i r_i x_i = 0,
\end{align*}

\end{exercise}


\begin{solution}


\end{solution}



\begin{exercise}

Esempio di media pesata con valori non compatibili.

\end{exercise}


\begin{solution}


\end{solution}



\begin{exercise}

Si generi un insieme di $N$ (ad esempio 10000) numeri pseudo-casuali distribuiti
secondo una distribuzione gaussiana in forma standard, ovvero con media $\mu = 0$
e deviazione standard $\sigma = 1$. Si usino i numeri per riempire un
istrogramma con un numero di canali $n$ fissato (ad esempio $n = 100$ tra
$-5$ e $+5$) e si faccia un fit con un modello gaussiano, stimando i valori di
\foreign{best-fit} dei parametri, le incertezze associate, il valore del
$\chisq$ ed il \pvalue.

Come variano tutte queste grandezze al variare di $n$ (ad esempio facendo variare
$n$ tra $5$ e $5000$ su una griglia equispaziata logaritmicamente)? Si spieghi in
modo quantitativo quello che sta succedendo.

\end{exercise}


\begin{solution}


\end{solution}



\begin{exercise}

Supponiamo di avere due grandezze fisiche $x$ e $y$ legate tra di loro da una
relazione funzionale di tipo lineare
\begin{align*}
  y = m x + q,
\end{align*}
dove $m$ e $q$ sono due parametri incogniti.

Se abbiamo a disposizione $10$ misure e possiamo scegliere liberamente i valori
di $x$ a cui effettuarle, come li sceglieremmo per minimizzare l'incertezza
su
\begin{enumerate}
    \item l'intercetta $q$?
    \item il coefficiente angolare $m$?
\end{enumerate}

\end{exercise}


\begin{solution}


\end{solution}



\begin{exercise}

Supponiamo di voler stimare il pi\`u precisamente possibile il tempo di decadimento
$\tau$ di un certo isotopo radioattivo a partire da un numero fissato (e.g., $10$)
misure indipendenti del rate \foreign{R} di decadimento, assumendo un andamento
esponenziale nel tempo
\begin{align*}
  R(t) = R_0 e^{-t/\tau}.
\end{align*}

Se siamo liberi di scegliere arbitrariamente i tempi $t_i$ a cui effettuare le
misure, qual \`e la strategia ottimale da adottare per minimizzare l'incertezza
sul valore di $\tau$? (Si assuma, per comodit\`a, di disporre di una stima
grossolana iniziale di $\tau$ per pianificare la misura.)

\end{exercise}


\begin{solution}


\end{solution}